\let\negmedspace\undefined
\let\negthickspace\undefined
\documentclass[journal,12pt,onecolumn]{IEEEtran}
\usepackage{cite}
\usepackage{amsmath,amssymb,amsfonts,amsthm}
\usepackage{algorithmic}
\usepackage{graphicx}
\usepackage{textcomp}
\usepackage{xcolor}
\usepackage{txfonts}
\usepackage{listings}
\usepackage{enumitem}
\usepackage{mathtools}
\usepackage{gensymb}
\usepackage{comment}
\usepackage[breaklinks=true]{hyperref}
\usepackage{tkz-euclide} 
\usepackage{gvv}                                        
%\def\inputGnumericTable{}                                 
\usepackage[latin1]{inputenc}     
\usepackage{xparse}
\usepackage{color}                                            
\usepackage{array}                                            
\usepackage{longtable}                                       
\usepackage{calc}                                             
\usepackage{multirow}
\usepackage{multicol}
\usepackage{hhline}                                           
\usepackage{ifthen}                                           
\usepackage{lscape}
\usepackage{tabularx}
\usepackage{array}
\usepackage{float}
\newtheorem{theorem}{Theorem}[section]
\newtheorem{problem}{Problem}
\newtheorem{proposition}{Proposition}[section]
\newtheorem{lemma}{Lemma}[section]
\newtheorem{corollary}[theorem]{Corollary}
\newtheorem{example}{Example}[section]
\newtheorem{definition}[problem]{Definition}
\newcommand{\BEQA}{\begin{eqnarray}}
\newcommand{\EEQA}{\end{eqnarray}}
\newcommand{\define}{\stackrel{\triangle}{=}}
\theoremstyle{remark}
\newtheorem{rem}{Remark}
% Marks the beginning of the document
\begin{document}
\title{2.10.11}
\author{ai25btech11022 - Narshitha}
\maketitle
\renewcommand{\thefigure}{\theenumi}
\renewcommand{\thetable}{\theenumi}
\textbf{Question}:\\
A unit vector that is coplanar with $\hat{i}+\hat{j}+2\hat{k}$ and $\hat{i}+2\hat{j}+\hat{k}$ and perpendicular to $\hat{i}+\hat{j}+\hat{k}$ is\_\_\_\_\\
\textbf{Solution}:\\
Let the two vectors lie on plane $\vec{P}$\\
as origin also lies on the plane\\
The Equation of the plane is $\vec{r}.{\vec{n}^T}=0$
\begin{align}
    \myvec{1\\1\\2}.{\vec{n}}^T=0\\
    \myvec{1\\2\\1}.{\vec{n}}^T=0\\
    \myvec{1 && 1 && 2\\1 && 2 &&1}.\vec{n}=\myvec{0\\0}\\
\end{align}
The augmented  matrix is
\begin{align}
    \myvec{1 && 1 && 2 &&\vrule&& 0\\1 && 2 && 1 &&\vrule&& 0}\\
\end{align}  
$R_2\rightarrow R_2-R_1$ and $R_1\rightarrow R_1-R_2$\\
\begin{align}
    \myvec{1 && 0 && 3 && \vrule&& 0\\0 && 1 && -1 &&\vrule&& 0}
\end{align}
we get\\
\begin{align}
    \vec{n}=\myvec{-3\\1\\1}
\end{align}
Let the unit vector be $\vec{A}$\\

\begin{align}
\vec{B}=\myvec{1\\1\\1}\\
{\myvec{\vec{n} && \vec{B}}}^T.\vec{A}=\myvec{0\\0}\\
    \myvec{-3 && 1 && 1\\1 && 1 && 1}.\vec{A}=\myvec{0\\0}
\end{align}
The augmented matrix is
\begin{align}
    \myvec{-3 && 1 && 1 && \vrule &&0\\ 1 && 1 && 1 &&\vrule&& 0}
\end{align}
$R_1\longrightarrow R_1-R_2$ and $R_2 \longrightarrow R_2 + R_1 /4$
\begin{align}
    \myvec{-4 && 0 && 0 && \vrule &&0\\0&& 1 && 1 && \vrule&&0}\\
    \vec{A}=\myvec{0\\b\\-b}
\end{align}
As $\vec{A}$ is a unit vector
\begin{align}
 \therefore  \vec{A}=\myvec{0\\\frac{1}{\sqrt{2}}\\-\frac{1}{\sqrt{2}}}
\end{align}
\begin{figure}[H]
    \centering
    \includegraphics[width=0.5\linewidth]{figs/fig1.png}
    \caption{ }
    \label{fig:placeholder}
\end{figure}
\end{document}
